\documentclass[11pt]{amsart}
\usepackage{amsmath,amssymb,amsthm,mathtools}
\usepackage[margin=1in]{geometry}
\usepackage{hyperref}

\newtheorem{theorem}{Theorem}[section]
\newtheorem{proposition}[theorem]{Proposition}
\newtheorem{lemma}[theorem]{Lemma}
\newtheorem{corollary}[theorem]{Corollary}
\theoremstyle{definition}
\newtheorem{definition}[theorem]{Definition}
\theoremstyle{remark}
\newtheorem{remark}[theorem]{Remark}
\newtheorem{observation}[theorem]{Numerical observation}

\newcommand{\fa}{\mathfrak a}
\newcommand{\fp}{\mathfrak p}
\newcommand{\ft}{\mathfrak t}
\newcommand{\RH}{\mathrm{RH}}

\title[The Pontryagin-Index Route]{The off-line zeros as a Pontryagin index: a conditional reduction and numerical observations on the prime tail}
\author{David Alejandro Trejo Pizzo}
\thanks{Independent Researcher. \texttt{dtrejopizzo@gmail.com}}


\begin{document}
\begin{abstract}
We study the Weil form $\ft=\fa-\fp$ (archimedean minus prime) as a quadratic
form and track not its sign but its \emph{negative index} $\kappa$, the number
of its negative squares. Weil's criterion gives $\kappa=0$ under RH; the
intermediate condition $\kappa<\infty$ (finite defect, a Pontryagin space
$\Pi_\kappa$) is the setting of Kre\u\i n--Langer theory. We prove a
conditional reduction: if the archimedean form has compact resolvent on the
ambient Hilbert space and the prime tail beyond a cutoff $X$ is
relatively bounded by the archimedean form with bound $<1$, then $\kappa$ is
finite, with an explicit bound. On the space $\mathcal H_d$ of functions
supported in $[-d,d]$ the hypotheses hold, so $\kappa_d<\infty$ for every $d$,
and the open question is the uniformity in $d$. We also describe, qualitatively,
exploratory computations on finite-dimensional band-limited subspaces which
suggest that the form of the hypothesis with $C=0$ and constants uniform in the
band fails: the cutoff at which the prime tail becomes subordinate appeared to
grow with the band, and the number of near-critical directions of the
prime-versus-archimedean pencil appeared to grow with the band. These
computations are not reproduced here and do not decide whether $\kappa$ is
finite. No claim about RH is made.
\end{abstract}
\maketitle

\section{Introduction}
The Weil explicit formula \cite{Weil} realizes a quadratic functional
$\ft(g)=\sum_\rho h_g(\gamma_\rho)$ over the nontrivial zeros
$\tfrac12+i\gamma_\rho$, with $h_g(r)=\widehat g(r)\overline{\widehat g(\bar r)}$
the entire extension of $|\widehat g|^2$ (for compactly supported $g$), where
$\widehat g(r)=\int_{\mathbb R}g(x)e^{irx}\,dx$;
on-line zeros contribute nonnegatively, off-line zeros contribute indefinite
quartets, and $\RH\iff\ft\succeq0$ (Weil's criterion \cite{Weil}). On
band-limited tests the zero side can be written, with $F=\widehat g$ and $E$
the evaluation map at the zeros, $\ft(F)=\langle JEF,EF\rangle$
with a fundamental symmetry $J$ whose negative part is supported on the
off-line zeros \cite[Theorem~1.2]{P8}. We write
\[
\kappa\;=\;\#\{\text{negative squares of }\ft\}\in\{0,1,2,\dots,\infty\}
\]
for its negative index, the maximal dimension of a subspace on which $\ft$ is
negative definite. We take $\kappa$ on compactly supported $g$, so
$\kappa=\sup_d\kappa_d$, where $\kappa_d$ is the negative index of $\ft$ on
$\mathcal H_d=\{g\in L^2(\mathbb R):\operatorname{supp}g\subset[-d,d]\}$.
Heuristically, by the identity just quoted, since each off-line quartet
contributes an indefinite term, one expects $\kappa$ to be finite exactly when only finitely
many zeros lie off the line, and then to be comparable with their number; we
do not prove this and use $\kappa$ only as defined above.

This paper asks whether $\kappa$ is \emph{finite}. An affirmative
answer is implied by RH ($\kappa=0$). If $\kappa<\infty$ and $\ft$ is
nondegenerate, completion gives a Pontryagin space $\Pi_\kappa$, in which every
self-adjoint operator has at most $\kappa$ eigenvalues in the open upper
half-plane (Pontryagin's theorem; see \cite{KL} and the references there). Finite defect is an
infinite-dimensional analogue of the finite-rank models that the function-field
proof of RH supplies geometrically.

We prove a conditional reduction of $\kappa<\infty$ to a relative form bound
on the prime tail (\S\ref{sec:red}), show that its hypotheses hold on each $\mathcal H_d$, and describe
exploratory computations (\S\S\ref{sec:c1}--\ref{sec:c2}) suggesting that a
$d$-uniform bound (RFB$_X$) with $C=0$ fails.

\section{The reduction to a tail relative form bound}\label{sec:red}
\subsection*{Setting}
For $g\in L^2(\mathbb R)$ write $\tilde g(x)=\overline{g(-x)}$ and
$(g\star\tilde g)(s)=\int_{\mathbb R}g(x+s)\overline{g(x)}\,dx=\langle\tau_{-s}g,g\rangle$,
where $\tau_s g(x)=g(x-s)$ and $\|g\|$ is the $L^2(\mathbb R)$ norm. The prime
form is
\[
\fp[g]=2\sum_{n\ge2}\frac{\Lambda(n)}{\sqrt n}\,\mathrm{Re}\,(g\star\tilde g)(\log n),
\]
and the archimedean form is
\[
\fa[g]=\frac1{2\pi}\int_{\mathbb R}\Omega(r)\,|\widehat g(r)|^2\,dr
+h_g(\tfrac i2)+h_g(-\tfrac i2),
\qquad \Omega(r)=\mathrm{Re}\,\psi(\tfrac14+\tfrac{ir}{2})-\log\pi,
\]
the last two terms being the contributions of the poles $s=0,1$ in the explicit
formula \cite{Weil} (with $\widehat g$ as above, $h_g(\pm\tfrac i2)
=\widehat g(\pm\tfrac i2)\overline{\widehat g(\mp\tfrac i2)}$, and
$\widehat g(\pm\tfrac i2)=\int g(x)e^{\mp x/2}dx$). For $g\in\mathcal H_d$
these terms are finite, $(g\star\tilde g)(s)=0$ for $|s|>2d$, and so the prime
sum is finite (only $n\le e^{2d}$ occur). We work on a closed subspace
$\mathcal H\subset L^2(\mathbb R)$ of test functions on which these
expressions are defined (for instance $\mathcal H_d$), and we define
$\ft=\fa-\fp$ on the domain of $\fa$; for smooth $g\in\mathcal H_d$ this agrees
with the sum over zeros by the explicit formula \cite{Weil}. Split the prime form at a cutoff
$X$, $\fp=\fp_{\le X}+\fp_{>X}$ (terms $n\le X$ and $n>X$). Since
$|(g\star\tilde g)(s)|\le\|g\|^2$, the head is bounded:
\begin{equation}\label{eq:head}
|\fp_{\le X}[g]|\le B_X\|g\|^2,\qquad B_X=2\sum_{n\le X}\Lambda(n)n^{-1/2}.
\end{equation}
It is not of finite rank: each term is the form of $(\tau_s+\tau_{-s})/2$,
which has continuous spectrum on $L^2(\mathbb R)$.

\begin{proposition}[Reduction]\label{prop:red}
Assume that $\fa$ is a closed form on $\mathcal H$, bounded below, whose
associated self-adjoint operator $A$ has compact resolvent (its eigenvalues
tend to $+\infty$). If for some $X$ the tail satisfies the one-sided relative
form bound
\begin{equation}\tag{RFB$_X$}
\fp_{>X}[g]\le \alpha\,\fa[g]+C\|g\|^2\qquad(g\in\operatorname{dom}\fa),\qquad 0\le\alpha<1,
\end{equation}
then $\ft$ has finitely many negative squares on $\operatorname{dom}\fa$, and
\[
\kappa\le N:=\#\{\text{eigenvalues }\mu\text{ of }A\text{ with }(1-\alpha)\mu<C+B_X\},
\]
counted with multiplicity. In particular (RFB$_X$) with the two-sided bound
$|\fp_{>X}[g]|\le\alpha\fa[g]+C\|g\|^2$ suffices.
\end{proposition}
\begin{proof}
By (RFB$_X$) and \eqref{eq:head},
$\ft[g]\ge(1-\alpha)\fa[g]-(C+B_X)\|g\|^2$ for $g\in\operatorname{dom}\fa$.
Let $E$ be the closed span of the eigenvectors of $A$ with
$(1-\alpha)\mu<C+B_X$; $N=\dim E<\infty$ by compactness of the resolvent. For
$g\in\operatorname{dom}\fa\cap E^\perp$, the spectral theorem gives
$\fa[g]\ge\mu_\ast\|g\|^2$ with $(1-\alpha)\mu_\ast\ge C+B_X$, hence
$\ft[g]\ge0$. A subspace on which $\ft$ is negative definite meets
$E^\perp$ only in $0$, so its dimension is at most $N$.
\end{proof}

\begin{remark}
On $\mathcal H=L^2(\mathbb R)$ the hypothesis cannot hold: the pole terms are
not defined there, and the multiplier $\Omega$ has continuous spectrum. On $\mathcal H_d$ it holds. Indeed $\Omega$ is
continuous, $\Omega(r)=\log|r|+O(1)$ as $|r|\to\infty$ (Stirling), hence
$\Omega\ge-c$. The form $g\mapsto\frac1{2\pi}\int(\Omega+c)|\widehat g|^2$ is
closed and nonnegative on $L^2(\mathbb R)$ (a multiplication form), so its
restriction to the closed subspace $\mathcal H_d$ is closed; its domain contains
$C_c^\infty(-d,d)$, hence is dense in $\mathcal H_d$. The pole terms satisfy
$|\widehat g(\pm\tfrac i2)|\le e^{d/2}\sqrt{2d}\,\|g\|$ on $\mathcal H_d$, so they
form a bounded form of rank at most $2$, and $\fa$ is closed and bounded below
on $\mathcal H_d$. Choose $c'$ with $\fa[g]+c'\|g\|^2\ge0$ on $\mathcal H_d$. For
compactness of the resolvent it suffices that the set
$K=\{g\in\mathcal H_d:\fa[g]+(c'+1)\|g\|^2\le1\}$, the unit ball of the form
domain, is precompact in $L^2$. Functions in $K$ are supported in $[-d,d]$ and
satisfy $\int_{|r|>R}|\widehat g|^2\le C_d/\inf_{|r|>R}(\Omega+c)$, with
$C_d$ depending only on $d$, $c$ and $c'$; the right-hand side tends to $0$
as $R\to\infty$, uniformly on $K$. By the Kolmogorov--Riesz criterion (uniform
localisation in $x$ and in $r$) $K$ is precompact. Moreover $\fp$ involves only
$n\le e^{2d}$ on $\mathcal H_d$. Proposition~\ref{prop:red} with
$X=e^{2d}$, $\alpha=0$, $C=0$ then gives
$\kappa_d\le\#\{\text{eigenvalues of }A_d\text{ below }B_{e^{2d}}\}<\infty$
for every $d$, where $A_d$ is the operator of $\fa$ on $\mathcal H_d$. The open
question is whether $\kappa_d$ stays bounded as $d\to\infty$; (RFB$_X$) for a
fixed $X$ with constants independent of $d$ bounds $\kappa_d$ by the number of
eigenvalues $\mu$ of $A_d$ with $(1-\alpha)\mu<C+B_X$, a number which itself may
grow with $d$.
\end{remark}

For a finite-dimensional subspace $V\subset\mathcal H_d$ on which $\fa\succ0$
(in \S\S\ref{sec:c1}--\ref{sec:c2}, a space spanned by finitely many functions
whose Fourier transforms are sinc functions at spacing $\pi/d$ centred near a
frequency $T_0$; we call such a space a band-$d$ space at height $T_0$), define
\begin{equation}\label{eq:lmax}
\lambda_{\max}(X,d)=\sup_{0\ne g\in V}\frac{\fp_{>X}[g]}{\fa[g]}.
\end{equation}
If (RFB$_X$) holds with $C=0$ on $\mathcal H_d$, then
$\lambda_{\max}(X,d)\le\alpha$ for every such $V\subset\mathcal H_d$;
conversely, $\lambda_{\max}(X,d)\le\alpha$ gives (RFB$_X$) with $C=0$ on $V$
only. Both statements are immediate from \eqref{eq:lmax}. The quantity
$\lambda_{\max}(X,d)$ depends on $V$, that is, on $T_0$ and on the dimension,
which the notation suppresses. On a single band the tail is a finite sum, so the content of
(RFB$_X$) is the uniformity in $d$ for fixed $X$.

\begin{remark}[Why this is hard]
$\sum_{n\le N}\Lambda(n)n^{-1/2}\sim2\sqrt N$ dominates the logarithmic growth
of $\Omega$, so a termwise estimate cannot give (RFB$_X$); heuristically it can hold
only through cancellation in the prime tail.
\end{remark}

\section{Computation I: the subordination cutoff grows with the
band}\label{sec:c1}
We discretised the forms on band-$d$ spaces at height $T_0=10^3$ where
$\fa\succ0$, and computed $\lambda_{\max}(X,d)$ as the largest generalized
eigenvalue of the pencil $(\fp_{>X},\fa)$. These were exploratory computations;
their code and data are not part of this paper and the basis dimensions are not
recorded, so we state only the qualitative behaviour observed, as numerical
observations that the reader cannot check from this paper.

\begin{observation}\label{obs:c1}
For $d=4,5$, the computed $\lambda_{\max}(1,d)$ was equal to $1$ within the
working precision, and the cutoff $X^\ast(d)$ at which $\lambda_{\max}(X,d)$
first dropped below $1$ was a fixed fraction of $e^{2d}$ (of order $e^{2d}/10$),
so that it grew with the band.
\end{observation}

\begin{observation}\label{cor:c1}
If the behaviour of Observation~\ref{obs:c1} persists, then for every fixed $X$
the bound (RFB$_X$) with $C=0$ cannot hold uniformly in $d$. This does not
decide whether $\kappa$ is finite: (RFB$_X$) is only a sufficient condition,
and the case $C>0$ is not tested.
\end{observation}

\section{Computation II: near-critical directions of the full
pencil}\label{sec:c2}
We next asked whether the value $\lambda_{\max}\approx1$ of the full pencil
$(\fp,\fa)$ is carried by few directions, computing its spectrum on band-$d$
spaces for $d=3,\dots,6$ (same discretisation and same caveat as in
\S\ref{sec:c1}).

\begin{observation}\label{obs:c2}
At $T_0=10^3$, the number of generalized eigenvalues of $(\fp,\fa)$ close to
$1$ increased with $d$ over $d=3,\dots,6$, roughly linearly, and the top
eigenvalues were nearly degenerate.
\end{observation}

\begin{observation}\label{cor:c2}
In these computations, a subspace of growing dimension satisfies
$\fp\approx\fa$, i.e.\ $\ft\approx0$. If the growth persists, the near-null
space of $\ft$ is not finite-dimensional in the limit $d\to\infty$, and no
finite critical subspace carries the near-saturation. This does not decide
whether $\kappa$ is finite.
\end{observation}

\section{Conclusion}
Proposition~\ref{prop:red} reduces finite defect, under a compact-resolvent
hypothesis on $\fa$, to a one-sided relative bound on the prime tail; on each
$\mathcal H_d$ it gives $\kappa_d<\infty$. The exploratory computations of
\S\S\ref{sec:c1}--\ref{sec:c2} suggest that a $d$-uniform bound (RFB$_X$)
with $C=0$ fails: in the computed range the prime tail is not made subordinate
to $\fa$ (with $C=0$) by a cutoff $X$ independent of the band, and the
near-saturation
of the pencil $(\fp,\fa)$ is spread over a subspace whose dimension grows with
the band. As a heuristic remark, this linear growth in the band is reminiscent
of the $\log T$ growth of the zero density, but we have no computation
supporting a precise relation. Nothing here decides whether $\kappa$ is finite,
and no claim about RH is made.

\begin{thebibliography}{9}
\bibitem{KL} M.~G.~Kre\u\i n and H.~Langer, \emph{\"Uber die
$Q$-Funktion eines $\pi$-hermiteschen Operators im Raume $\Pi_\kappa$}, Acta
Sci. Math. (Szeged) \textbf{34} (1973), 191--230.
\bibitem{P8} D.~A.~Trejo Pizzo, \emph{A Kre\u\i n-space formulation of the zero
side of the Weil quadratic form on Paley--Wiener spaces}, preprint (Paper 08).
\bibitem{Weil} A.~Weil, \emph{Sur les ``formules explicites'' de la th\'eorie
des nombres premiers}, Comm.\ S\'em.\ Math.\ Univ.\ Lund (1952), 252--265.
\end{thebibliography}

\end{document}
